\subsection{داده، تقسیم استاندارد و ممیزی منشأ}

هر سطر داده یک تصویر است، اما تصویرهای یک بیمار یا یک معاینه مشاهده‌های مستقل نیستند؛ همین وابستگی مبنای ساخت تقسیم استاندارد قرار گرفت. جمعیت کانونی گزارش ۸۸۶۲ مورد کامل است و ترکیب آن به تفکیک منبع و کلاس در جدول \ref{tab:meth_cohort} آمده است.

\subsubsection{جمعیت کانونی و منابع داده}

داده از سه منبع تصویر فوندوس ساخته شد. \lr{Plus}\enfoot{نام یکی از منابع تصویربرداری و نیز برچسب مرحله‌ی پلاس در رتینوپاتی نارس؛ معادل انگلیسی: \lr{Plus}} بزرگ‌ترین منبع است و فراداده‌ی دستگاه \lr{RetCam}\enfoot{دستگاه تصویربرداری شبکیه‌ی نوزادان؛ معادل انگلیسی: \lr{RetCam}} را همراه دارد. \lr{FARFUM-RoP}\enfoot{مجموعه‌ی عمومی تصاویر رتینوپاتی نارس؛ معادل انگلیسی: \lr{FARFUM-RoP}} و داده‌ی فارابی دو منبع دیگرند. برچسب‌های \lr{FARFUM-RoP} از پرونده‌ی صفحه‌گسترده خوانده شد و در فارابی شناسه‌ی بیمار موجود نبود، پس شناسه‌ی یکتای معاینه جای آن را گرفت.

هشت مورد که در سه نشانگر فراکتالی مقدار \lr{NaN}\enfoot{مقدار عددی تعریف‌نشده؛ معادل انگلیسی: \lr{NaN}} داشتند کنار گذاشته شدند و همه‌ی آن‌ها در بخش آموزش بودند. هیچ جایگزینی برای مقدارهای گمشده انجام نشد؛ بنابراین جمعیت گزارش، موارد کامل است. سهم هر منبع و هر کلاس در جدول \ref{tab:meth_cohort} و شکل‌های شکل مربوط و شکل مربوط آمده است.

\begin{table}[htbp]
  \centering
  \small
  \caption{ترکیب جمعیت نهایی موارد کامل به تفکیک منبع تصویربرداری}
  \label{tab:meth_cohort}
  \begin{tabular}{lrrrrrr}
    \toprule
    منبع & تصویر & گروه & \lr{Normal} & \lr{Pre-Plus} & \lr{Plus} & درصد \\
    \midrule
    فارابی & ۱۴۰۹ & ۱۶۰ & ۳۶۹ & ۴۵۲ & ۵۸۸ & \lr{۱۵.۹} \\
    \lr{FARFUM-RoP} & ۱۵۲۸ & ۶۸ & ۷۸۰ & ۴۷۸ & ۲۷۰ & \lr{۱۷.۲} \\
    \lr{Plus} & ۵۹۲۵ & ۱۸۶ & ۵۳۰۴ & ۰ & ۶۲۱ & \lr{۶۶.۹} \\
    \midrule
    مجموع & ۸۸۶۲ & ۴۱۴ & ۶۴۵۳ & ۹۳۰ & ۱۴۷۹ & \lr{۱۰۰.۰} \\
    \bottomrule
  \end{tabular}
\end{table}

\begin{figure}[htbp]
  \centering
  \imgbox[1.0\textwidth]{figures/report/main/D1_source_class_distribution.pdf}
  \caption{توزیع کلاس در جمعیت نهایی به تفکیک منبع تصویربرداری (8862 تصویر، 414 گروه). منبع \lr{Plus} هیچ مورد \lr{Pre-Plus} ندارد.}
  \label{fig:U1}
\end{figure}

\begin{figure}[htbp]
  \centering
  \imgbox[1.0\textwidth]{figures/report/main/D3_source_population.pdf}
  \caption{سهم هر منبع در جمعیت نهایی.}
  \label{fig:U2}
\end{figure}

\begin{figure}[htbp]
  \centering
  \imgbox[1.0\textwidth]{figures/report/main/D2_split_class_distribution.pdf}
  \caption{توزیع کلاس در تقسیم‌های آموزش، اعتبارسنجی و آزمون.}
  \label{fig:U3}
\end{figure}




\subsubsection{فضای برچسب}

سه برچسب \lr{Normal}\enfoot{وضعیت طبیعی شبکیه؛ معادل انگلیسی: \lr{Normal}}، \lr{Pre-Plus}\enfoot{وضعیت پیش از مرحله‌ی پلاس؛ معادل انگلیسی: \lr{Pre-Plus}} و \lr{Plus} به سه کد صفر، یک و دو نگاشته می‌شوند. منبع \lr{Plus} هیچ مورد \lr{Pre-Plus} ندارد؛ پس برای این منبع نه مدل سه‌کلاسه معنا دارد و نه \lr{AUC}\enfoot{سطح زیر منحنی مشخصه‌ی عملکرد گیرنده؛ معادل انگلیسی: \lr{AUC}} سه‌کلاسه. در آن منبع فقط معیار محدود «\lr{Normal} در برابر \lr{Plus}» با برچسب صریح گزارش می‌شود و هرگز به‌جای مقدار سه‌کلاسه نمی‌نشیند. این تفکیک لازم است، زیرا ممکن است مدلی دو گروه منفی را با هم اشتباه بگیرد و همچنان \lr{Plus} را درست رتبه‌بندی کند.

\subsubsection{تقسیم استاندارد}

شمارش هر بخش و هر کلاس در جدول \ref{tab:meth_split} و شکل شکل مربوط آمده است. تخصیص گروه‌ها به سه بخش با یک رویه‌ی قطعی انجام می‌شود؛ هیچ گروهی در بیش از یک بخش قرار نمی‌گیرد و اجرای دوباره‌ی همان رویه همان تقسیم را می‌سازد. بخش آموزش بیشترین سهم را از منبع \lr{Plus} دارد و فارابی کمترین سهم را.

\begin{table}[htbp]
  \centering
  \small
  \caption{تقسیم گروه‌محور استاندارد جمعیت موارد کامل}
  \label{tab:meth_split}
  \begin{tabular}{lrrrrr}
    \toprule
    تقسیم & \lr{N} & \lr{Normal} & \lr{Pre-Plus} & \lr{Plus} & گروه \\
    \midrule
    آموزش & ۶۲۰۳ & ۴۴۹۰ & ۶۵۰ & ۱۰۶۳ & ۲۹۰ \\
    اعتبارسنجی & ۱۳۲۸ & ۹۸۱ & ۱۴۰ & ۲۰۷ & ۶۲ \\
    آزمون & ۱۳۳۱ & ۹۸۲ & ۱۴۰ & ۲۰۹ & ۶۲ \\
    \midrule
    مجموع & ۸۸۶۲ & ۶۴۵۳ & ۹۳۰ & ۱۴۷۹ & ۴۱۴ \\
    \bottomrule
  \end{tabular}
\end{table}



تعریف گروه میان منابع یکسان نیست. در \lr{Plus} و \lr{FARFUM-RoP} گروه بر پایه‌ی بیمار ساخته شده است، اما در فارابی شناسه‌ی بیمار در دسترس نبود و گروه بر پایه‌ی معاینه تعریف شد. نتیجه‌ی مستقیم این تفاوت آن است که برای فارابی، نبود هم‌پوشانی در سطح معاینه تأیید شده است و استقلال در سطح بیمار اثبات‌نشده باقی می‌ماند؛ به همین دلیل در سراسر این گزارش فقط عبارت «نبود هم‌پوشانی در سطح شناسه‌ی گروه قابل‌دسترس» به‌کار می‌رود و هیچ ادعای مطلقی فراتر از همین سطح مطرح نمی‌شود.

\subsubsection{جایگاه آزمون استاندارد در برابر تحلیل‌های ثانویه}

تقسیم استاندارد پیش از این برای مقایسه‌ی اصلی به‌کار رفته و همه‌ی پارامترهای قابل انتخاب با بخش اعتبارسنجی تعیین شده‌اند. بنابراین هر تحلیل دیگری که روی همان آزمون اجرا شده باشد، تحلیل مجدد اکتشافی ثانویه روی یک تقسیم قفل‌شده است و آزمون تازه محسوب نمی‌شود. این تفکیک در سراسر گزارش حفظ می‌شود و مستندسازی این تحلیل‌ها به‌معنای تأیید دوباره‌ی آن‌ها نیست. هر عدد گزارش‌شده به همین تقسیم، به بذر پروژه و به اثر انگشت پیکربندی متناظر خود گره خورده است.
